\chapter{Machine Learning Pipelines with Azure ML and Databricks}
\index{Azure Machine Learning}\index{AML pipeline}\index{AutoML}\index{compute cluster}\index{managed feature store}\index{MLOps}%
\begin{objectives}
\begin{itemize}
    \item Navigate the Azure Machine Learning workspace: studio, compute instances, compute clusters, datastores, and environments.
    \item Choose among AutoML, the Designer, and custom notebooks for a given team and problem.
    \item Register ADLS Gen2 data via datastores secured with the workspace managed identity.
    \item Build an AML pipeline from versioned components---preprocess, train, evaluate, register---and execute it from the Python SDK.
    \item Register features in the managed feature store so training and serving share one definition.
    \item Integrate Databricks for heavy feature computation feeding AML training.
    \item Design a secure workflow for data scientists training PyTorch models on GPU clusters over governed lake data.
\end{itemize}
\end{objectives}

\begin{crosscloud}
Managed ML platforms share a universal anatomy: a workspace, versioned data references, reusable environments, pipeline components, and a model registry.

\vspace{2mm}
{\footnotesize
\begin{tabular}{@{}p{3.0cm}p{2.9cm}p{3.1cm}p{3.2cm}@{}}
\toprule
\textbf{Capability} & \textbf{AWS} & \textbf{Azure} & \textbf{GCP} \\
\midrule
ML platform & SageMaker & Azure Machine Learning & Vertex AI \\
AutoML & SageMaker Autopilot & Azure AutoML & Vertex AutoML \\
Pipelines & SageMaker Pipelines & AML pipelines (components) & Vertex Pipelines (Kubeflow) \\
Feature store & SageMaker Feature Store & AML managed feature store & Vertex AI Feature Store \\
Model registry & SageMaker Model Registry & AML Model Registry & Vertex Model Registry \\
Notebook compute & Studio notebooks & Compute instances & Workbench instances \\
\addlinespace
\multicolumn{4}{@{}l@{}}{\textbf{Fabric}: built-in ML with MLflow on lakehouses; \textbf{Snowflake}: Snowpark ML; \textbf{MongoDB}: partner MLOps over Atlas.} \\
\bottomrule
\end{tabular}}

\vspace{1mm}
MLflow is the lingua franca underneath: AML, Databricks, and Fabric all track experiments with it, so experiment-tracking skills transfer even where the platform does not.
\end{crosscloud}

\begin{keyterms}
\textbf{Workspace}: the AML management boundary---compute, data, experiments, models, endpoints.\quad
\textbf{Compute instance}: a managed single-node development VM; \textbf{compute cluster}: autoscaling multi-node training compute.\quad
\textbf{Environment}: a versioned, reproducible software context (Docker image plus conda/pip spec).\quad
\textbf{Component}: a reusable, versioned pipeline step with declared inputs, outputs, and environment.\quad
\textbf{Datastore}: a registered, credential-managed reference to storage such as ADLS Gen2.
\end{keyterms}

\section{Week 18 Overview: The ML Platform Layer}
Chapter 17 embedded a trained model into SQL. Chapters 18--19 build the platform that produces, ships, and watches such models continuously. Azure Machine Learning is that platform: a workspace unifying compute, data access, experiment tracking, pipelines, a model registry, and endpoints \cite{microsoft2025azureml}. Databricks remains in the picture where it belongs---heavy feature computation over the lake---with AML owning training orchestration, registration, and serving.

The discipline of the week is \emph{reproducibility through versioning}: versioned data assets, versioned environments, versioned components, versioned models. A pipeline built from versioned pieces can answer the only question that matters when a model misbehaves in production: \emph{exactly what produced this artifact?}

\section{Monday: Workspace, Compute, and Data Access}
\index{Azure Machine Learning!workspace}\index{compute cluster}

\subsection{Compute Instances vs.\ Compute Clusters}
The \textbf{compute instance} is a per-developer managed VM for notebooks---always-on, single-node, your desk in the cloud. The \textbf{compute cluster} is the training fleet: autoscaling, multi-node, GPU-capable, and---critically---\emph{able to scale to zero}. The cost rule mirrors Chapter 9's job clusters: development on instances, training on clusters, and a zero minimum node count so idle time bills nothing.

\subsection{Datastores and Managed Identity}
A datastore registers ADLS Gen2 (or other storage) in the workspace. Register it with the workspace's \textbf{managed identity} and grant that identity \texttt{Storage Blob Data Reader} on the lake: training jobs then read governed data with no keys, no service-principal secrets, and an auditable identity \cite{microsoft2025rbac}. Data assets---named, versioned references over datastore paths---give pipelines stable inputs (\texttt{azureml:features:3}) even as the underlying folders change.

\begin{pitfall}
\textbf{Notebook-mounted storage with personal credentials.} A datastore with a personal account key in a notebook bypasses every governance control the lake has and dies when the person changes roles. Managed-identity datastores are the only production pattern.
\end{pitfall}

\section{Tuesday: AutoML, Designer, or Custom Code}
\index{AutoML}\index{Azure ML Designer}

\subsection{The Three Authoring Modes}
\textbf{AutoML} sweeps algorithms and hyperparameters against a dataset and metric, producing a ranked leaderboard with explainability---the right tool for baselines and for teams establishing ``how hard is this problem, really.'' The \textbf{Designer} composes drag-and-drop pipelines from modules---accessible, but its artifacts are the least code-reviewable of the three. \textbf{Custom notebooks and components} give full control (PyTorch, custom losses, exotic features) at full responsibility. The mature pattern uses AutoML to set the baseline bar and custom components to beat it; a custom model that cannot beat AutoML's leaderboard entry has not earned its complexity.

\subsection{The Thursday Preview: AutoML to Endpoint in One Afternoon}
The fastest path through the platform: point AutoML at a registered data asset, let it sweep, deploy the best run to a \textbf{managed online endpoint}, and test with a REST payload. This is Week 18's original project, and it remains the fastest way to establish the end-to-end skeleton (data asset $\rightarrow$ trained model $\rightarrow$ registered model $\rightarrow$ endpoint) before replacing the middle with your own pipeline.

\section{Wednesday: AML Pipelines---Components and Environments}
\index{AML pipeline}\index{component}\index{environment}

\subsection{Pipelines from Versioned Parts}
An AML pipeline is a DAG of \textbf{components} \cite{microsoft2025amlPipelines}. Each component declares inputs, outputs, a command, and an \textbf{environment}; each is versioned independently and reused across pipelines. Environments pin the software world (base image plus dependencies) so a pipeline that ran in March reproduces in October. The canonical pipeline shape is Preprocess $\rightarrow$ Train $\rightarrow$ Evaluate $\rightarrow$ (conditionally) Register.

\begin{lstlisting}[language=Python,caption={An AML pipeline assembled from registered components via the v2 SDK.},label={lst:ch18-pipeline}]
from azure.ai.ml import dsl, Input
from azure.ai.ml.entities import Environment

env = Environment(
    name="churn-env", version=3,
    image="mcr.microsoft.com/azureml/openmpi4.1.0-ubuntu22.04",
    conda_file="conda.yml",
)

@dsl.pipeline()
def churn_pipeline(raw: Input):
    prep  = preprocess_component(input_data=raw)
    train = train_component(train_data=prep.outputs.clean)
    eval_ = evaluate_component(model=train.outputs.model,
                               test_data=prep.outputs.test)
    return {"model": train.outputs.model}

ml_client.jobs.create_or_update(
    churn_pipeline(raw=Input(path="azureml:features:1")))
\end{lstlisting}

\subsection{The Managed Feature Store}
AML's managed feature store operationalizes Chapter 17's contract: feature sets are versioned, materialized to offline (lake) and online (cache) stores, and carry point-in-time-correct retrieval for training-set assembly. Features stop being notebook-local code and become cataloged, monitored assets with owners---the ML counterpart of the medallion gold zone.

\subsection{Databricks in the Picture}
Where feature computation is heavy---sessionization over billions of events, graph features---Databricks does the work on the lake and writes the feature tables; AML pipelines then consume them as data assets. The boundary rule: Databricks owns \emph{feature materialization}, AML owns \emph{training orchestration and the registry}; the feature table in the governed lake is the contract between them.

\section{Thursday Hands-on Project: Preprocess, Train, Evaluate, Register}
\index{hands-on lab}\index{AML pipeline!project}
Build the canonical AML pipeline from the Python SDK and execute it from a notebook: preprocess the churn features, train, evaluate, and register the model only if it clears a metric gate.

\subsection{Provision the Workspace}
\begin{lstlisting}[language=bash,caption={Workspace with a managed identity and a scale-to-zero training cluster.},label={lst:ch18-aml}]
$rg = "rg-de-azure-book-week18"
az group create --name $rg --location eastus
az extension add --name ml
az ml workspace create --name aml-week18 --resource-group $rg --location eastus

az ml compute create --name cpu-cluster --resource-group $rg `
  --workspace-name aml-week18 --type AmlCompute `
  --size Standard_DS3_v2 --min-instances 0 --max-instances 4 `
  --idle-time-before-scale-down 300
\end{lstlisting}

Register the ADLS datastore with the workspace managed identity and grant it read access to the feature container.

\subsection{Author and Run the Pipeline}
Define \texttt{preprocess}, \texttt{train}, and \texttt{evaluate} components (each a small script plus environment), assemble them as in Listing~\ref{lst:ch18-pipeline}, and submit. In the studio, inspect the pipeline graph, per-step logs, and the tracked metrics.

\subsection{Gate Registration on a Metric}
Extend the evaluate component to emit \texttt{auc}, and add a condition: register only when $\text{AUC} \ge 0.80$. Run once with a deliberately degraded feature set (drop the strongest feature) and confirm the gate blocks registration---the failing case is the deliverable that proves the gate exists.

\section{Friday Use Case: GPU PyTorch over a Governed Lake}
\index{use case}\index{GPU}\index{PyTorch}
A media company's data scientists must train custom PyTorch recommendation models on GPU clusters against customer interaction data in the governed ADLS lake---without copying data to unmanaged storage and without any secrets in notebooks.

\subsection{Design}
Workspace managed identity holds \texttt{Storage Blob Data Reader} on the lake; a data asset versions each training snapshot. GPU compute clusters (NC-family) scale from zero per job. Environments pin CUDA/PyTorch versions; MLflow tracking captures every run; the registry holds promoted models with lineage back to the exact data asset, environment, and code version. Purview (Chapter 22) sees the lake side; AML's lineage covers the model side. The security review can trace every model to the data and identity that produced it---which is the actual requirement behind ``secure access to governed data.''

\begin{pitfall}
\textbf{Export-first workflows.} ``Export the lake to CSV for training'' breaks lineage, residency, and freshness in one step. The datastore pattern exists precisely so governed data is read \emph{in place} by an auditable identity.
\end{pitfall}

\section{Architectural Decision Record Template}
\index{architectural decision record}
Per pipeline record: component inventory with owners, environment pinning policy, data asset versions consumed, the metric gate and its threshold rationale, compute sizing (SKU, min/max nodes, idle scale-down), and the Databricks/AML responsibility boundary.

\section{Operational Deep Dive: Pipeline Schedules and Reuse}
\index{pipeline schedule}
AML pipelines run on schedules or via Event Grid triggers (data-arrival-driven training). Because components are versioned and cached, unchanged steps reuse prior outputs---re-running a pipeline after a code change to \texttt{train} does not re-run \texttt{preprocess}. Monitor pipeline duration per step; the step that doubled is usually the one whose data doubled, which is a feature-table question, not an ML question.

\section{Troubleshooting Patterns}
\index{troubleshooting!AML}
Job fails at environment build $\rightarrow$ dependency conflict in the conda spec; pin versions. Job fails on data access $\rightarrow$ the workspace identity lacks the lake RBAC role or the datastore fell back to key auth. Gate never passes after a data change $\rightarrow$ inspect the feature table before questioning the threshold. Cluster never scales $\rightarrow$ max instances or regional quota; check both before resizing the job.

\section{Week 18 Lab Rubric}
\index{rubric}
\begin{table}[H]
\centering
\small
\begin{tabular}{@{}p{3.2cm}p{5.0cm}p{5.0cm}@{}}
\toprule
\textbf{Criterion} & \textbf{Meets expectation} & \textbf{Needs improvement} \\
\midrule
Reproducibility & Versioned data, environment, components; pipeline re-runs identically & Local notebooks and unpinned dependencies \\
Security & Managed-identity datastore; no secrets in notebooks & Personal keys or mounted storage \\
Gating & Metric gate blocks a demonstrated failing model & Every run registers a model \\
Cost & Cluster scales to zero; idle time bounded & Always-on GPU cluster \\
\bottomrule
\end{tabular}
\caption{Week 18 rubric for the AML pipeline deliverables.}
\label{tab:ch18-rubric}
\end{table}

\section{Interview Q\&A}
\subsection{Concepts and Trade-offs}
\begin{enumerate}
    \item \textbf{Q: When would you choose AutoML over a custom training notebook?}\\
    A: For baselines and well-structured tabular problems---AutoML sets the bar a custom model must beat to justify its complexity.
    \item \textbf{Q: Why use AML compute clusters instead of compute instances for training?}\\
    A: Clusters autoscale across nodes for training and scale to zero when idle; instances are single-node development desks.
    \item \textbf{Q: How does AML securely access governed ADLS training data?}\\
    A: Through a datastore bound to the workspace managed identity with lake RBAC---no keys, full audit trail.
    \item \textbf{Q: What does model registration give you beyond saving a file?}\\
    A: Versioned identity with lineage to data, environment, code, and metrics---the answer to ``what exactly produced this artifact.''
\end{enumerate}

\section{Further Reading}
The Azure ML overview \cite{microsoft2025azureml}, the pipelines concept guide \cite{microsoft2025amlPipelines}, and the examples repository \cite{microsoft2025azuremlExamples} are the primary sources. Chapter 19 continues with endpoints, drift, and the retraining loop; Chapter 17 supplies the feature contract this pipeline consumes.

\section*{Key Takeaways}
\begin{itemize}
    \item Instances for development, clusters for training, minimum nodes zero: ML compute economics mirror Spark's.
    \item Managed-identity datastores are the only governed path from the lake to training.
    \item AutoML sets the baseline; custom components must beat it to justify their complexity.
    \item Versioned components plus pinned environments make pipelines reproducible and cache-efficient.
    \item The managed feature store turns features into cataloged, point-in-time-correct assets.
    \item Registration is lineage: every production model answers ``what data, code, and environment made me.''
\end{itemize}

\section{Exercises}
\begin{enumerate}
    \item \textbf{(Conceptual)} Explain why a zero-minimum-node cluster changes team behavior, not just cost.
    \item \textbf{(Conceptual)} What breaks when two pipelines share a component version and one team edits it?
    \item \textbf{(Conceptual)} Why does the Designer produce less reviewable artifacts than components-as-code?
    \item \textbf{(Hands-on)} Complete the Thursday project, including the gated-registration failing case.
    \item \textbf{(Hands-on)} Run an AutoML sweep on the same data asset and compare its best AUC against your pipeline's model.
    \item \textbf{(Hands-on)} Change only the \texttt{train} component and observe step-level caching on re-run; record which steps executed.
    \item \textbf{(Design)} Design the Databricks/AML boundary for the media company: which features are materialized where, and why.
    \item \textbf{(Design)} Write the ADR for the metric gate: threshold, metric choice, and the override process.
    \item \textbf{(Design)} Specify an Event Grid trigger that retrains when the feature table's drift metric crosses a threshold (preview Chapter 19).
    \item \textbf{(Interview)} Prepare a two-minute answer to ``why did this model change behavior after a re-run with no code changes?''
\end{enumerate}

\section{Capstone Project: Gated Training Pipeline over Governed Features}\label{sec:ch18-capstone}
\index{capstone project}\index{AML pipeline!capstone}

\begin{capstonebox}
\textbf{Mission:} Deliver a reproducible training pipeline: versioned data assets read via managed identity, pinned environments, componentized preprocess/train/evaluate steps, a metric-gated registration, and AutoML as the recorded baseline.

\textbf{Estimated time:} 5--6 hours.

\textbf{What you will practice:}
\begin{itemize}
    \item Workspace, compute, and datastore provisioning with identity-based access.
    \item Component authoring and pipeline assembly with the v2 SDK.
    \item Metric gates with a proven failing case.
    \item Baseline discipline: AutoML versus custom, on record.
\end{itemize}
\end{capstonebox}

\subsection*{Scenario}
The subscription business graduates from Chapter 17's hand-run training to a pipeline the platform team can schedule. Data science owns components; the platform team owns the gate; compliance requires full lineage from model to training data.

\subsection*{Requirements}
\begin{itemize}
    \item \textbf{Access:} all data read through the managed-identity datastore; zero secrets in code.
    \item \textbf{Reproducibility:} pinned environment; versioned components and data assets; re-run produces identical metrics.
    \item \textbf{Gate:} registration only at $\text{AUC} \ge 0.80$, with the failing case demonstrated.
    \item \textbf{Baseline:} an AutoML run on the same asset, with the comparison recorded in the model description.
\end{itemize}

\subsection*{Reference Architecture}
\begin{figure}[H]
    \centering
    \resizebox{\textwidth}{!}{%
    \begin{tikzpicture}[node distance=1.2cm]
        \node[store] (lake) {ADLS features\\(datastore, MI)};
        \node[stage, right=1.3cm of lake] (pipe) {AML pipeline\\prep $\rightarrow$ train $\rightarrow$ eval};
        \node[stage, right=1.3cm of pipe] (gate) {Metric gate\\AUC $\ge$ 0.80};
        \node[store, right=1.3cm of gate] (reg) {Model\\Registry};
        \node[tool, above=0.9cm of pipe] (automl) {AutoML\\baseline};
        \node[doc, below=0.9cm of pipe] (cluster) {Compute cluster\\min 0 / max 4};
        \draw[arrow] (lake) -- (pipe);
        \draw[arrow] (pipe) -- (gate);
        \draw[arrow] (gate) -- (reg);
        \draw[arrow, dashed] (automl) -- (gate);
        \draw[arrow, dashed] (cluster) -- (pipe);
    \end{tikzpicture}%
    }
    \caption{Capstone architecture: governed features feed a gated pipeline; AutoML provides the recorded baseline.}
    \label{fig:ch18-capstone-arch}
\end{figure}

\subsection*{Implementation Steps}
\begin{enumerate}
    \item Provision workspace and compute (Listing~\ref{lst:ch18-aml}); register the datastore and data asset.
    \item Author the three components and assemble the pipeline (Listing~\ref{lst:ch18-pipeline}).
    \item Add the metric gate; demonstrate pass and fail runs.
    \item Run the AutoML baseline; record the comparison in the registry description.
    \item Verify lineage: registry entry resolves to data asset, environment, code, and metrics.
\end{enumerate}

\subsection*{Deliverables}
\begin{itemize}
    \item Component and pipeline code, environment specs, and provisioning scripts.
    \item Gate pass/fail evidence and the AutoML comparison record.
    \item A README that reproduces the pipeline from a clean workspace.
\end{itemize}

\subsection*{Acceptance Criteria}
\begin{itemize}
    \item A re-run with unchanged inputs reuses cached steps and reproduces metrics.
    \item A degraded feature set is blocked from registration by the gate.
    \item Every registered model traces to its data asset version and environment.
\end{itemize}

\subsection*{Stretch Goals}
\begin{itemize}
    \item Add the managed feature store and swap the pipeline input to a versioned feature set.
    \item Schedule the pipeline and add step-duration regression alerting.
    \item Move feature materialization to a Databricks job and consume its output as the data asset.
\end{itemize}
